\section{从关中到建康：统一是一串接收动作}
\label{sec:sui-curated-chronicle}

开皇元年二月，北周静帝的禅让仪式在长安完成；仪式之外，关中军政集团、北方边防和南方的陈朝仍是三种不同的现实。\eventanchor{ST-581-SUI-ACCESSION}{隋受禅（581）}杨坚得到的是可供调动的中枢，不是一张已经服从的疆域图。北方突厥压力使朝廷不能把南征当作唯一事务；对旧周官僚、州县和户籍的接收，也要先在纸面、仓廪和道路上成为日常工作。%
\footnote{受禅年序见《隋书·高祖纪上》《北史·隋本纪上》与《资治通鉴》开皇元年二月条。三者足以互校改朝的顺序，不能据仪文推定宫廷成员的真实意图。}

开皇八年，水陆诸军南下；次年正月建康陷落，陈后主被俘。\eventanchor{ST-588-SOUTHERN-CAMPAIGN}{隋军南下（588）}\eventanchor{ST-589-CHEN-CONQUEST}{建康易手（589）}长江防线被穿透，依赖的是北方整合、边防暂得缓冲和陈朝军政困境同时出现，而不是一个可以抽离粮道、渡口与地方接应的“北强南弱”故事。胜者本纪与《陈书》的亡国叙事都留下了宫廷画面；它们能确认行动的先后，不能把兵数、路线细节或江南人的态度写成确数。

\begin{turningpoint}{建康失守不是江南已经被接管}
城门的转换结束了陈朝的中央政权，却开始了新朝在江南登记户口、安置军民和重置州县的难题。开皇十年的江南起事\eventanchor{ST-590-JIANGNAN-REBELLION}{江南起事（590）}提醒我们：军事胜利与地方服从之间隔着差役、官吏、交通和暴力。
\end{turningpoint}

大业初年，运河工程把洛阳、江淮与北方军需接入更密的水路网络；它给粮运、人员与命令提供了新的通道，也把征发和地方负担推到更多家庭面前。\eventanchor{ST-605-TONGJI-CANAL}{开凿通济渠（605）}随后对高句丽的远征把边疆战争、转运与征役叠在一起。\eventanchor{ST-612-GOGURYEO-INVASION}{高句丽远征（612）}这里不宜用一条工程或一场战争解释隋亡；可见的是中央动员的规模扩大，地方承受与边地战略却未必同步。

\begin{sourcenote}{运河、远征与民变：材料能说到何处}
《隋书·炀帝纪》《隋书·食货志》与《资治通鉴》保存了工程和战争的年序；唐人修史常把它们编入亡国因果。条目、里程和军数适于比较制度语言，不足以量出每一县的徭役、死亡或反抗动机。敦煌、墓志与城址可补局部尺度，却不能替代全域账目。
\end{sourcenote}

隋的遗产由此是双面的：统一的行政尺度、交通和法律成为唐可继承的资源，而以远距离动员解决边防和财政问题的压力也留了下来。读到太原起兵时，可转向\eventref{ST-617-LI-YUAN-TAIYUAN}{太原起兵（617）}，看另一支关中军政网络如何接住这套未完成的设施与名义。
